\documentclass[11pt,a4paper]{article}
\usepackage[utf8]{inputenc}
\usepackage[T1]{fontenc}
\usepackage[magyar]{babel}
\usepackage{amsmath,amssymb}
\usepackage{geometry}
\geometry{top=25mm, bottom=25mm, left=25mm, right=25mm}
\usepackage{tcolorbox}
\usepackage{xcolor}
\usepackage{tikz}

\tcbuselibrary{skins,breakable}

\definecolor{myblue}{RGB}{26, 82, 118}
\definecolor{mybg}{RGB}{242, 247, 250}

\newtcolorbox{defbox}[2][]{
  colback=mybg,
  colframe=myblue,
  fonttitle=\bfseries,
  title=#2,
  rounded corners,
  boxrule=0.8pt,
  #1
}

\newtcolorbox{examplebox}[1][]{
  colback=white,
  colframe=teal!80!black,
  fonttitle=\bfseries,
  title=Megoldott példa: #1,
  rounded corners,
  boxrule=0.8pt
}

\newtcolorbox{exercisebox}{
  colback=orange!5!white,
  colframe=orange!80!black,
  fonttitle=\bfseries,
  title=Önálló feladatok,
  rounded corners,
  boxrule=0.8pt
}

\title{Deriválás összefoglaló}
\author{Matematika Összefoglaló}
\date{}

%%% Megjelenítéshez automatikusan kiegészített csomagok %%%
\usepackage{amsfonts, amsthm}
\usepackage{pgfplots}
\pgfplotsset{compat=1.18}
\usepackage{graphicx}
\usepackage{enumitem}
\usepackage{booktabs}
\usepackage{float}
\usepackage{caption}
\usepackage{microtype}

\usepackage{hyperref}

\begin{document}

\maketitle

\section{Bevezetés és a derivált fogalma}
A differenciálszámítás a matematikai elemzés egyik alappillére. Geometriailag a derivált egy görbe érintőjének meredekségét, fizikailag pedig a pillanatnyi sebességet vagy változási ütemet jelenti.

\begin{defbox}{A derivált definíciója}
Legyen az $f$ függvény értelmezve az $x_0$ pont egy környezetében. Azt mondjuk, hogy az $f$ függvény differenciálható az $x_0$ helyen, ha létezik a következő véges határérték:
\[
f'(x_0) = \lim_{x \to x_0} \frac{f(x) - f(x_0)}{x - x_0} = \lim_{h \to 0} \frac{f(x_0 + h) - f(x_0)}{h}
\]
Ezt a határértéket az $f$ függvény $x_0$ pontbeli deriváltjának nevezzük.
\end{defbox}

\section{Deriválási szabályok}
Ha $f$ és $g$ differenciálható függvények egy adott intervallumon, és $c$ valós szám, akkor a következő alapvető műveleti szabályok érvényesek:

\begin{itemize}
    \item \textbf{Konstans szorzó:} $(c \cdot f(x))' = c \cdot f'(x)$
    \item \textbf{Összeg és különbség:} $(f(x) \pm g(x))' = f'(x) \pm g'(x)$
    \item \textbf{Szorzatfüggvény deriváltja:} $(f(x) \cdot g(x))' = f'(x)g(x) + f(x)g'(x)$
    \item \textbf{Hányadosfüggvény deriváltja:} $\left(\frac{f(x)}{g(x)}\right)' = \frac{f'(x)g(x) - f(x)g'(x)}{(g(x))^2}$, feltéve, hogy $g(x) \neq 0$.
\end{itemize}

\section{Elemi függvények deriváltjai}
Az alábbi listában összefoglaljuk a legfontosabb elemi függvények deriváltjait:

\begin{align*}
    (c)' &= 0 \quad (\text{ahogyan } c \text{ konstans}) \\
    (x^n)' &= n \cdot x^{n-1} \quad (n \in \mathbb{R}) \\
    (e^x)' &= e^x, \quad (a^x)' = a^x \ln a \\
    (\ln x)' &= \frac{1}{x}, \quad (\log_a x)' = \frac{1}{x \ln a} \\
    (\sin x)' &= \cos x, \quad (\cos x)' = -\sin x \\
    (\operatorname{tg} x)' &= \frac{1}{\cos^2 x}, \quad (\operatorname{ctg} x)' = -\frac{1}{\sin^2 x}
\end{align*}

\section{Láncszabály és összetett függvények deriválása}
Ha egy függvény egy másik függvénybe van ágyazva (összetett függvény), akkor a deriváláshoz a láncszabályt használjuk.

\begin{defbox}{A láncszabály}
Ha az $f$ függvény differenciálható a $g(x)$ helyen, és a $g$ függvény differenciálható az $x$ helyen, akkor az $F(x) = f(g(x))$ összetett függvény is differenciálható, és deriváltja:
\[
F'(x) = f'(g(x)) \cdot g'(x)
\]
Szavakkal: a külső függvény deriváltja a belső függvényben kiszámolva szorozva a belső függvény deriváltjával.
\end{defbox}

\begin{examplebox}{Összetett függvény deriválása}
Határozzuk meg az $f(x) = \sin(3x^2 + 1)$ függvény deriváltját!
\vspace{0.2cm}

\textbf{Megoldás:}
Itt a külső függvény a szinusz ($\sin(u)$), a belső függvény pedig $g(x) = 3x^2 + 1$.
A külső függvény deriváltja: $\cos(u)$, így $\cos(3x^2 + 1)$.
A belső függvény deriváltja: $g'(x) = 6x$.
A láncszabály szerint:
\[
f'(x) = \cos(3x^2 + 1) \cdot 6x = 6x \cos(3x^2 + 1)
\]
\end{examplebox}

\section{Inverz függvények deriválása és az arkuszfüggvények}
Ha egy $f$ függvénynek létezik differenciálható inverze ($f^{-1}$), és $f'(f^{-1}(y)) \neq 0$, akkor az inverz függvény is deriválható, és teljesül a következő összefüggés:
\[
(f^{-1})'(y) = \frac{1}{f'(f^{-1}(y))}
\]

Ennek a képletnek a segítségével származtathatók az inverz trigonometrikus függvények (arkuszfüggvények) deriváltjai:

\begin{align*}
    (\arcsin x)' &= \frac{1}{\sqrt{1 - x^2}} \quad (-1 < x < 1) \\
    (\arccos x)' &= -\frac{1}{\sqrt{1 - x^2}} \quad (-1 < x < 1) \\
    (\operatorname{arctg} x)' &= \frac{1}{1 + x^2} \quad (x \in \mathbb{R}) \\
    (\operatorname{arcctg} x)' &= -\frac{1}{1 + x^2} \quad (x \in \mathbb{R})
\end{align*}

\section{Hiperbolikus függvények deriváltja}
A hiperbolikus függvények exponenciális kifejezésekkel definialt függvények, amelyek nagyon hasonlítanak a trigonometrikus függvényekre, de tulajdonságaikban eltérnek.

\begin{defbox}{Hiperbolikus alapfüggvények}
\[
\operatorname{sh} x = \frac{e^x - e^{-x}}{2}, \quad \operatorname{ch} x = \frac{e^x + e^{-x}}{2}, \quad \operatorname{th} x = \frac{\operatorname{sh} x}{\operatorname{ch} x}
\]
\end{defbox}

Ezen függvények deriváltjai rendkívül egyszerűek, nincs bennük előjélváltás a szinuszhoz hasonlóan:

\begin{align*}
    (\operatorname{sh} x)' &= \operatorname{ch} x \\
    (\operatorname{ch} x)' &= \operatorname{sh} x \\
    (\operatorname{th} x)' &= \frac{1}{\operatorname{ch}^2 x}
\end{align*}

\begin{examplebox}{Inverz hiperbolikus / arkusz kombináció}
Számítsuk ki a $g(x) = \operatorname{arctg}(\operatorname{sh} x)$ függvény deriváltját!
\vspace{0.2cm}

\textbf{Megoldás:}
Alkalmazzuk a láncszabályt. A külső függvény az $\operatorname{arctg}$, aminek a deriváltja $\frac{1}{1 + u^2}$, a belső függvény pedig $\operatorname{sh} x$.
\[
g'(x) = \frac{1}{1 + (\operatorname{sh} x)^2} \cdot \operatorname{sh}'(x) = \frac{\operatorname{ch} x}{1 + \operatorname{sh}^2 x}
\]
Mivel az alapvető hiperbolikus összefüggés szerint $\operatorname{ch}^2 x - \operatorname{sh}^2 x = 1$, azaz $1 + \operatorname{sh}^2 x = \operatorname{ch}^2 x$:
\[
g'(x) = \frac{\operatorname{ch} x}{\operatorname{ch}^2 x} = \frac{1}{\operatorname{ch} x}
\]
\end{examplebox}

\section{Önálló feladatok}
Gyakorold a fent leírt szabályokat az alábbi feladatokon keresztül!

\begin{exercisebox}
Határozd meg az alábbi függvények első deriváltját:
\begin{enumerate}
    \item $f(x) = 5x^4 - 3\ln x + 2e^x$
    \item $f(x) = (x^2 + 1) \cdot \operatorname{tg} x$
    \item $f(x) = e^{-2x} \cdot \sin(5x)$
    \item $f(x) = \arcsin(2x)$
    \item $f(x) = \operatorname{ch}(3x^2 - 1)$
\end{enumerate}
\end{exercisebox}

\end{document}